\chapter{Hernia Repair}

\section*{Introduction \& Clinical Indications}
Hernia repair, medically referred to as herniorrhaphy or hernioplasty, is a prevalent surgical intervention aimed at correcting a hernia, which is characterized by the protrusion of an organ or fatty tissue through a defect in the surrounding muscle or connective tissue. This procedure is most commonly performed on the abdominal wall, addressing various types of hernias, including inguinal, femoral, umbilical, and incisional hernias. The clinical significance of hernia repair lies in its ability to restore anatomical integrity, alleviate symptoms such as pain and discomfort, and prevent serious complications like incarceration or strangulation, which can lead to ischemia and necrosis of the affected tissue. By reinforcing the weakened area, hernia repair not only improves the patient's quality of life but also reduces the risk of recurrence and associated morbidity.

\begin{itemize}
  \item Herniorrhaphy
  \item Hernioplasty
  \item Inguinal Hernia Repair
  \item Ventral Hernia Repair
  \item Laparoscopic Hernia Repair
  \item Open Hernia Repair
  \item Groin Hernia Repair
\end{itemize}

The primary surgical principle of hernia repair is to reduce the herniated contents back into their anatomical cavity and to close or reinforce the fascial defect through which they protruded. This is achieved through either primary tissue repair, where the edges of the defect are sutured together, or by utilizing prosthetic reinforcement, typically with synthetic mesh. The latter, known as tension-free repair, is the preferred method for most adult hernias due to its significantly lower recurrence rates compared to primary tissue repairs, which often create tension on the suture line, leading to tissue ischemia and eventual failure.

The rationale for surgical intervention centers on restoring the integrity of the abdominal wall, preventing further herniation, and alleviating symptoms such as pain, discomfort, or a visible bulge. For inguinal hernias, the goal is to repair the defect in the transversalis fascia or posterior wall of the inguinal canal. For ventral and incisional hernias, the aim is to close the fascial gap in the anterior abdominal wall, often employing component separation techniques for larger defects. Mesh repair functions by providing a scaffold for tissue ingrowth, creating a strong, durable barrier that resists future herniation. Technical goals include minimizing tissue dissection to preserve neurovascular structures, achieving a tension-free repair to reduce recurrence, and ensuring adequate mesh overlap beyond the defect margins to provide robust reinforcement.

The history of hernia repair is rich and spans centuries, with early attempts involving rudimentary methods such as bandaging and cautery. Significant advancements began in the late 19th century, driven by improved anatomical understanding and the introduction of antiseptic techniques. In 1884, Edoardo Bassini, an Italian surgeon, made a landmark contribution by describing a technique that involved high ligation of the hernia sac and reconstruction of the posterior wall of the inguinal canal by suturing the conjoint tendon to the inguinal ligament. This "Bassini repair" dramatically reduced recurrence rates from previous methods, which often exceeded 30-50\%, to around 10-15\%, establishing it as the gold standard for decades.

Further evolution in hernia repair techniques was influenced by William Halsted in the United States, who modified Bassini's approach by excising the cremaster muscle and transposing the spermatic cord. The mid-20th century saw the introduction of prosthetic materials, with Francis Usher pioneering the use of polypropylene mesh for hernia repair in 1958, which paved the way for tension-free repairs. Irving Lichtenstein popularized the "Lichtenstein tension-free hernioplasty" in the 1980s, utilizing a flat piece of mesh placed over the defect without tension on the suture line, which became widely adopted due to its simplicity, reproducibility, and significantly lower recurrence rates, typically less than 5\%.

The late 20th century marked the advent of minimally invasive techniques, with Ralph Ger performing the first laparoscopic hernia repair in 1982, involving mesh plug repair. Subsequent refinements by pioneers such as Barry McKernan and George Ferzli led to the development of laparoscopic transabdominal preperitoneal (TAPP) and totally extraperitoneal (TEP) repairs, which offer reduced postoperative pain, faster recovery, and comparable recurrence rates to open mesh repairs. The latest evolution in hernia repair is robotic-assisted techniques, which provide enhanced dexterity and three-dimensional visualization.

Hernia repair is indicated for a variety of clinical scenarios, primarily aimed at alleviating symptoms and preventing complications. The following are clear clinical indications for surgical intervention:

\begin{itemize}
  \item To alleviate pain or discomfort caused by the hernia, which may worsen with activity, coughing, or straining.
  \item To prevent or treat incarceration, where the herniated tissue becomes trapped and cannot be manually reduced, leading to persistent pain and swelling.
  \item To prevent or treat strangulation, a life-threatening complication where the blood supply to the incarcerated tissue is cut off, leading to ischemia, necrosis, and potential sepsis.
  \item For hernias that are progressively increasing in size, posing a higher risk of future complications and making repair more challenging.
  \item To address cosmetic concerns related to a visible bulge, particularly in the abdominal wall, which can impact body image and quality of life.
  \item When conservative management, such as watchful waiting or truss use, is deemed inappropriate or has failed to control symptoms or prevent progression.
  \item For femoral hernias, which have a significantly higher propensity for incarceration and strangulation compared to inguinal hernias, often warranting earlier repair even if asymptomatic.
  \item In cases of recurrent hernias following previous repair attempts, where a different surgical approach or technique may be required.
  \item For pediatric hernias (e.g., inguinal, umbilical) that persist beyond infancy or are symptomatic, as spontaneous closure is less likely.
\end{itemize}

Hernia repair is predominantly a primary treatment for symptomatic hernias. For patients experiencing pain, discomfort, or a progressively enlarging hernia, elective surgical repair is the standard of care. This is considered the first-line intervention to restore anatomical integrity and prevent future complications. In cases of asymptomatic inguinal hernias, particularly in men, a period of watchful waiting may be considered; however, surgical repair remains the definitive treatment if symptoms develop or if the patient prefers intervention due to lifestyle concerns or anxiety about potential complications.

Conversely, hernia repair becomes an urgent or emergency procedure when complications arise. Specifically, for incarcerated hernias that cannot be reduced, or critically, for strangulated hernias where the blood supply to the herniated contents is compromised, immediate surgical intervention is required. In these emergency scenarios, the procedure is life-saving, acting as a salvage operation to prevent tissue necrosis, sepsis, and potentially death. Thus, while often elective, hernia repair serves as a critical emergency procedure in specific, high-risk scenarios, highlighting its dual role in clinical practice.

\section*{Relevant Surgical Anatomy}
A thorough understanding of the relevant surgical anatomy is essential for the successful execution of hernia repair. The abdominal wall is composed of several layers, including the skin, subcutaneous tissue, fascia, and muscles, which provide structural support. 

For **inguinal hernias**, the inguinal canal is a critical anatomical feature, measuring approximately 4 cm in length and extending from the deep inguinal ring (an opening in the transversalis fascia) to the superficial inguinal ring (an opening in the external oblique aponeurosis). The canal contains the spermatic cord in males, which includes the vas deferens, testicular artery, pampiniform plexus, and the genitofemoral nerve. In females, the round ligament occupies this space. The posterior wall of the inguinal canal is primarily formed by the transversalis fascia, and direct inguinal hernias protrude through Hesselbach's triangle, bordered by the inferior epigastric vessels laterally, the rectus abdominis medially, and the inguinal ligament inferiorly. Indirect inguinal hernias follow the path of the spermatic cord through the deep inguinal ring.

For **femoral hernias**, the defect occurs through the femoral canal, located medial to the femoral vein and artery, just inferior to the inguinal ligament. The femoral ring, which is the superior opening of the canal, is bounded by the inguinal ligament anteriorly, the lacunar ligament medially, the superior pubic ramus posteriorly, and the femoral vein laterally. Femoral hernias are particularly concerning due to their narrow confines, which predispose them to incarceration and strangulation.

**Umbilical hernias** arise from a defect in the linea alba at the umbilical ring, where the umbilical vessels once traversed. **Ventral and incisional hernias** occur through defects in the anterior abdominal wall fascia, often at sites of previous surgical incisions or in the linea alba. The rectus abdominis muscles, along with their fascial sheaths and the oblique muscles, contribute to the integrity of the abdominal wall. The arterial supply to the abdominal wall is primarily provided by the superior and inferior epigastric arteries, while venous drainage corresponds to these vessels. Lymphatic drainage in the groin region is primarily through the superficial and deep inguinal lymph nodes.

\section*{Preoperative Workup \& Preparation}
Thorough preoperative preparation is crucial for optimizing patient safety and surgical outcomes. This typically involves a comprehensive assessment to identify and mitigate risks.

\begin{itemize}
  \item \textbf{Medical Evaluation:} A comprehensive history and physical examination are performed to assess the patient's overall health, identify comorbidities (e.g., cardiac disease, diabetes, pulmonary issues), and review all current medications.
  
  \item \textbf{Laboratory Tests:} Routine blood tests, such as a complete blood count, electrolyte panel, renal function tests, and coagulation studies, are usually ordered. Urinalysis may also be performed to rule out urinary tract infection.
  
  \item \textbf{Imaging Studies:} While not always necessary for straightforward clinical hernias, imaging like ultrasound, CT scan, or MRI may be used to confirm the diagnosis, delineate complex anatomy, assess for complications (e.g., incarceration), or plan for recurrent or large incisional hernias.
  
  \item \textbf{Cardiac and Anesthesia Clearance:} Patients with significant cardiac or pulmonary conditions (e.g., history of myocardial infarction, severe COPD) may require preoperative cardiac evaluation (e.g., EKG, stress test, echocardiogram) and clearance from an anesthesiologist to assess fitness for general anesthesia.
  
  \item \textbf{Medication Adjustments:} Patients are instructed to discontinue blood-thinning medications (e.g., aspirin, clopidogrel, warfarin, novel oral anticoagulants) several days to a week prior to surgery, as directed by the surgeon and cardiologist, to minimize bleeding risk. Other medications, such as insulin, may require dosage adjustments.
  
  \item \textbf{NPO Status:} Patients are typically instructed to refrain from eating or drinking for 6-8 hours before surgery to minimize the risk of aspiration during anesthesia induction.
  
  \item \textbf{Antibiotic Prophylaxis:} A single dose of intravenous prophylactic antibiotics (e.g., cefazolin) is often administered shortly before the incision to reduce the risk of surgical site infection, particularly when prosthetic mesh is used.
  
  \item \textbf{Informed Consent:} The patient must fully understand the procedure, its potential risks, expected benefits, and available alternatives, and provide written informed consent.
  
  \item \textbf{Bowel Preparation:} For very large ventral or incisional hernias, or those involving potential bowel resection (e.g., strangulated hernia), a mechanical bowel preparation and oral antibiotics may be required.
\end{itemize}

\textbf{Absolute Contraindications:}
\begin{itemize}
  \item Active systemic infection or sepsis, which significantly increases the risk of surgical site infection, mesh contamination, and systemic complications.
  
  \item Uncontrolled coagulopathy or severe bleeding disorders that cannot be corrected preoperatively, posing an unacceptable risk of hemorrhage.
  
  \item Severe, uncompensated medical conditions (e.g., acute myocardial infarction, unstable angina, severe heart failure, acute respiratory distress syndrome, uncontrolled diabetic ketoacidosis) that make the risks of general anesthesia and surgery prohibitive.
  
  \item Patient refusal or inability to provide informed consent.
\end{itemize}

\textbf{Relative Contraindications and Precautions:}
\begin{itemize}
  \item Active infection at the surgical site or in the skin overlying the hernia, which should be treated and resolved before elective repair to prevent mesh infection.
  
  \item Pregnancy, where elective repair is usually deferred until after delivery unless the hernia is acutely symptomatic, incarcerated, or strangulated.
  
  \item Morbid obesity (BMI > 40 kg/m\textasciicircum{}2), which increases technical difficulty, operative time, and risks of complications like wound infection, seroma, and recurrence. Weight loss is often recommended prior to elective repair.
  
  \item Severe chronic obstructive pulmonary disease (COPD) or other significant pulmonary disease, requiring careful anesthetic management, smoking cessation, and optimization of respiratory function preoperatively.
  
  \item Immunosuppression (e.g., due to chemotherapy, organ transplant, chronic corticosteroid use, HIV/AIDS), which increases infection risk and may impair wound healing.
  
  \item Uncontrolled diabetes mellitus (HbA1c > 8\%), which can significantly increase infection risk, impair wound healing, and contribute to cardiovascular complications.
  
  \item Patients on chronic anticoagulation or antiplatelet therapy, requiring careful management of these medications perioperatively, often involving bridging therapy or temporary cessation.
  
  \item Large, irreducible hernias with potential for "loss of domain" (where a significant portion of abdominal contents resides outside the abdominal cavity), requiring specialized preoperative planning, potentially including progressive pneumoperitoneum or botulinum toxin injections to relax abdominal wall muscles, and possibly staged repair.
\end{itemize}

\section*{Surgical Technique \& Procedure}
Hernia repair can be performed using various techniques, primarily open surgery or minimally invasive approaches (laparoscopic or robotic). The choice of technique depends on the type and size of the hernia, patient factors, surgeon expertise, and the presence of previous repairs.

The procedure typically begins with the administration of general anesthesia, although spinal or local anesthesia with sedation may be utilized for some open repairs. The patient is positioned supine, and the surgical site is prepped and draped in a sterile fashion.

\begin{itemize}
  \item \textbf{Open Inguinal Hernia Repair (e.g., Lichtenstein Tension-Free Hernioplasty):}
  \begin{enumerate}
    \item A 5-8 cm oblique incision is made in the groin crease, extending through the skin and subcutaneous tissue.
    \item The external oblique aponeurosis is incised, and its leaves are retracted to expose the inguinal canal.
    \item The spermatic cord (in males) is identified and carefully dissected from the hernia sac.
    \item The hernia sac is identified, separated from the cord structures, and its contents are gently reduced back into the abdominal cavity. The sac is then ligated at its neck and excised or inverted.
    \item A synthetic polypropylene mesh (typically 6x11 cm) is tailored and placed over the posterior wall defect, extending beyond the margins of Hesselbach's triangle.
    \item The mesh is secured to the pubic tubercle medially and the inguinal ligament inferiorly with non-absorbable sutures, ensuring a tension-free repair. The superior edge of the mesh is often sutured to the internal oblique muscle or conjoint tendon.
    \item The external oblique aponeurosis is reapproximated over the spermatic cord and mesh, and the subcutaneous tissue and skin are closed in layers.
  \end{enumerate}

  \item \textbf{Laparoscopic Inguinal Hernia Repair (e.g., TAPP - Transabdominal Preperitoneal):}
  \begin{enumerate}
    \item Three small keyhole incisions (5-12 mm) are made, typically one at the umbilicus for the camera port and two in the lower abdomen for working ports.
    \item Carbon dioxide gas is insufflated into the abdominal cavity to create pneumoperitoneum and working space.
    \item A laparoscope is inserted, and the hernia defect is visualized from within the abdomen.
    \item The peritoneum overlying the hernia defect is incised, and the herniated contents are reduced. A preperitoneal space is created by dissecting the peritoneum away from the abdominal wall.
    \item A large piece of synthetic mesh (typically 10x15 cm) is placed over the defect from the inside, extending well beyond its margins to cover all potential groin hernia sites (direct, indirect, femoral).
    \item The mesh is secured with absorbable tacks, sutures, or fibrin glue to the pubic bone, Cooper's ligament, and abdominal wall muscles.
    \item The incised peritoneum is then closed over the mesh, isolating it from the intra-abdominal contents.
    \item The small incisions are then closed with sutures or surgical glue.
  \end{enumerate}

  \item \textbf{Laparoscopic Totally Extraperitoneal (TEP) Repair:} This technique is similar to TAPP, but the entire procedure is performed in the preperitoneal space without entering the abdominal cavity, potentially reducing the risk of visceral injury.
\end{itemize}

Drains are rarely used for routine hernia repairs but may be considered in cases of large incisional hernias, significant dissection, or when a large seroma is anticipated.

\section*{Complications \& Risks}
While hernia repair is generally safe and effective, like any surgical procedure, it carries potential risks and complications. These can be broadly categorized as general surgical risks and procedure-specific risks.

\begin{itemize}
  \item \textbf{General Surgical Risks:}
  \begin{itemize}
    \item \textbf{Anesthesia Complications:} Adverse reactions to medications, respiratory problems (e.g., pneumonia, atelectasis), cardiac events (e.g., myocardial infarction, arrhythmia), or cerebrovascular events (stroke).
    
    \item \textbf{Bleeding:} Hematoma formation at the surgical site, which may require drainage, or rarely, significant intraoperative hemorrhage.
    
    \item \textbf{Infection:} Surgical site infection (SSI), which can involve the skin (cellulitis), deeper tissues, or the prosthetic mesh, potentially requiring prolonged antibiotics, wound debridement, or mesh removal.
    
    \item \textbf{Seroma:} Accumulation of clear, serous fluid at the surgical site, which may require aspiration, particularly common after large ventral hernia repairs.
    
    \item \textbf{Deep Vein Thrombosis (DVT) and Pulmonary Embolism (PE):} Blood clot formation in the deep veins of the legs that can dislodge and travel to the lungs, a potentially life-threatening complication.
  \end{itemize}

  \item \textbf{Procedure-Specific Risks:}
  \begin{itemize}
    \item \textbf{Hernia Recurrence:} The most common long-term complication, where the hernia reappears at the same or a nearby site, necessitating further surgery.
    
    \item \textbf{Chronic Postoperative Pain:} Persistent pain in the groin or abdominal wall lasting more than 3-6 months, often due to nerve entrapment, irritation, or mesh-related inflammation, particularly after open inguinal repair.
    
    \item \textbf{Nerve Injury:} Damage to specific nerves (e.g., ilioinguinal, iliohypogastric, genitofemoral nerves in inguinal repair) leading to numbness, burning, dysesthesia, or neuropathic pain in the distribution of the affected nerve.
    
    \item \textbf{Bowel or Visceral Injury:} Accidental damage to the intestine, bladder, blood vessels, or other abdominal organs, more common in laparoscopic or recurrent hernia repairs due to adhesions.
    
    \item \textbf{Urinary Retention:} Difficulty or inability to urinate after surgery, often requiring temporary catheterization, more common in elderly men with prostatic hypertrophy.
    
    \item \textbf{Testicular Complications (for inguinal hernia repair):} Ischemic orchitis (inflammation due to impaired blood supply to the testicle), testicular atrophy (shrinkage), or hydrocele formation (fluid collection around the testicle).
    
    \item \textbf{Mesh-Related Complications:} Mesh infection, erosion into adjacent organs (e.g., bowel, bladder), migration of the mesh, or chronic foreign body reaction leading to inflammation or fistula formation.
    
    \item \textbf{Adhesions:} Formation of scar tissue between organs and the mesh or surgical site, potentially leading to chronic pain or future bowel obstruction.
    
    \item \textbf{Vascular Injury:} Damage to major vessels such as the femoral artery or vein, or the inferior epigastric vessels, leading to significant bleeding or limb ischemia.
  \end{itemize}
\end{itemize}

\section*{Postoperative Care \& Recovery}
Immediately after hernia repair, the patient is transferred to the Post-Anesthesia Care Unit (PACU) for close monitoring as they recover from anesthesia. Vital signs, pain levels, and the surgical site are continuously assessed. Pain management is initiated, typically with oral or intravenous analgesics. Patients are encouraged to ambulate early, often within hours of surgery, to prevent complications like deep vein thrombosis, pulmonary embolism, and respiratory issues. Most elective hernia repairs, especially laparoscopic ones, are performed on an outpatient basis or with a single overnight hospital stay.

During the first 24-48 hours, patients may experience mild to moderate pain, swelling, and bruising at the surgical site. Oral pain medications (e.g., acetaminophen, NSAIDs, short-course opioids) are usually sufficient for discomfort. Activity is generally restricted to light ambulation, and patients are advised to avoid heavy lifting (typically >10-15 lbs), straining, or strenuous activities. Wound care instructions, including keeping the incision clean and dry, are provided. Over the first 1-2 weeks, pain gradually subsides, and patients can typically resume light daily activities, including driving (once off opioid pain medication). Stool softeners are often recommended to prevent constipation and straining during bowel movements, which could stress the repair. Full recovery and return to normal activities, including exercise and heavy lifting, usually take 4 to 6 weeks, depending on the size and type of hernia repair, the surgical approach, and the individual's healing capacity.

During hernia repair, the surgeon documents several key findings. These include the type of hernia (e.g., indirect, direct, femoral, umbilical, incisional), its size, whether it was reducible or incarcerated, and the contents of the hernia sac (e.g., omentum, small bowel, colon, bladder). The quality of the surrounding fascial tissue and the presence of any adhesions are also noted. For incarcerated or strangulated hernias, the viability of the herniated contents is critically assessed; non-viable tissue necessitates resection.

If a hernia sac or its contents are excised (e.g., a portion of omentum), these tissues are sent for pathological examination. The pathology report typically confirms the presence of adipose tissue (omentum), small or large bowel, or other visceral structures within the hernia sac. It may also describe inflammatory changes, signs of ischemia or necrosis if strangulation occurred, or the presence of specific tissue types. The report serves to confirm the nature of the herniated tissue and rule out any unexpected pathologies, although primary hernia repair rarely yields significant pathological findings beyond confirming the presence of normal tissue in an abnormal location.

A normal and successful postoperative course following hernia repair is characterized by several key indicators. Clinically, the patient should experience a progressive resolution of the preoperative pain and discomfort associated with the hernia. The abdominal wall or groin area should be flat, without any visible or palpable bulge at the site of the repair, indicating successful reduction and reinforcement. The surgical incision should show signs of healthy healing, with minimal redness, swelling, or discharge. Postoperative pain should be manageable with oral analgesics and progressively decrease over days to weeks, allowing for a gradual return to normal activities. There should be no signs of infection, significant hematoma, seroma requiring intervention, or new neurological deficits indicative of nerve injury. Ultimately, the goal is for the patient to return to their baseline functional status without restrictions or recurrence of the hernia, typically within 4-6 weeks.

Postoperative care is essential for ensuring proper healing and detecting any potential complications.

\begin{itemize}
  \item \textbf{Postoperative Clinic Visit:} Typically scheduled 1-2 weeks after surgery to assess wound healing, remove non-absorbable sutures or staples if present, and address any patient concerns regarding pain, activity, or symptoms.
  
  \item \textbf{Activity Resumption Guidance:} Patients receive specific instructions on a gradual increase in physical activity, with clear timelines for resuming light exercise, driving, and eventually heavy lifting or strenuous activities (usually 4-6 weeks post-op).
  
  \item \textbf{Pain Management Review:} Assessment of pain levels and adjustment of analgesics if necessary, with emphasis on transitioning from opioid to non-opioid pain relief.
  
  \item \textbf{Warning Signs Education:} Patients are thoroughly advised to contact their surgeon immediately if they experience fever (above 101.5\textdegree F or 38.6\textdegree C), worsening or severe pain, excessive swelling, spreading redness, pus or foul-smelling discharge from the incision, or a new bulge at the repair site.
  
  \item \textbf{Physical Therapy/Rehabilitation:} Rarely needed for routine hernia repair, but may be recommended for patients undergoing repair of very large incisional hernias, or those with significant deconditioning, to aid in core strengthening and functional recovery.
\end{itemize}

Long-term follow-up is generally not required unless symptoms or concerns about recurrence arise, at which point further evaluation may be warranted.

\section*{Outcomes \& Prognosis}
Hernias, particularly inguinal hernias, are among the most common surgical conditions worldwide, affecting millions of individuals annually. Inguinal hernias have a lifetime prevalence of approximately 27\% in men and 3\% in women, making them significantly more common in males. The incidence increases with age, peaking in the 7th and 8th decades of life, often associated with weakening of connective tissues. Umbilical hernias are also common, especially in infants (where many resolve spontaneously by age 4-5) and adults, with a higher prevalence in women and those with increased intra-abdominal pressure (e.g., pregnancy, obesity, ascites, chronic cough). Incisional hernias occur in approximately 10-15\% of patients following abdominal surgery, with risk factors including wound infection, obesity, poor surgical technique, and conditions that impair wound healing. Femoral hernias are less common, accounting for about 3-5\% of groin hernias, but have a higher risk of incarceration and strangulation, occurring more frequently in women. The overall mortality rate for elective hernia repair is very low, typically less than 0.1\%, but can rise significantly (up to 10-20\%) in cases of emergency repair for strangulated hernias, especially in elderly or comorbid patients. Morbidity rates vary by hernia type and technique, with chronic pain and recurrence being the most common complications.

The outcomes of modern hernia repair are generally excellent, with high success rates and significant improvement in patient quality of life. For elective inguinal hernia repair with mesh, success rates (defined as freedom from recurrence) typically range from 95-99\% at 5-10 years. Both laparoscopic (TAPP/TEP) and open (Lichtenstein) mesh repairs demonstrate comparable recurrence rates in experienced hands, with patient-reported outcomes often favoring laparoscopic approaches for reduced early postoperative pain and faster return to activity. The main long-term complication impacting quality of life is chronic postoperative pain, which occurs in 10-15\% of patients after inguinal hernia repair, though severe debilitating pain is less common (1-2\%).

Prognostic factors influencing outcomes include patient characteristics (e.g., obesity, smoking, connective tissue disorders, diabetes), hernia characteristics (size, type, recurrence status, presence of incarceration/strangulation), and surgical technique (tension-free repair, adequate mesh overlap, meticulous nerve identification and preservation). Recurrence rates are higher for large, recurrent, or complicated hernias, and in patients with risk factors that increase intra-abdominal pressure or impair wound healing. For example, the European Hernia Society guidelines emphasize the importance of mesh repair for adult inguinal hernias to minimize recurrence. Overall, the prognosis for patients undergoing elective hernia repair is very favorable, with a low risk of serious complications and a high likelihood of a durable repair and symptom resolution, allowing a full return to normal activities.

\section*{References}
\begin{enumerate}
  \item MedlinePlus. Hernia Repair. U.S. National Library of Medicine; 2024. Available from: \url{https://medlineplus.gov/herniarepair.html}
  
  \item Brunicardi FC, Andersen DK, Billiar TR, et al. Schwartz's Principles of Surgery. 11th ed. McGraw-Hill Education; 2019.
  
  \item Fitzgibbons RJ Jr, Forse RA. Clinical practice. Groin hernias in adults. N Engl J Med. 2015 Feb 19;372(8):756-63.
  
  \item European Hernia Society Guidelines on the Treatment of Inguinal Hernia in Adult Patients. Hernia. 2018 Feb;22(1):1-165.
\end{enumerate}

\clearpage
